\section{Related work}\label{sec:related}

\paragraph{NBA careers.} Player evaluation has moved from raw box-score counts to composite indices such as PER \citep{hollinger2002}, Win Shares, Box Plus-Minus and plus-minus ratings \citep{kubatko2007,hvattum2019,deshpande2016}, and to data-mining approaches that segment players by role rather than nominal position \citep{sarlis2020,anilduman2024,pereztoledano2019,radovanovic2013}. Longitudinal questions have mostly been framed as \emph{aging curves}: population-level models of performance against age, estimated with Bayesian hierarchical models \citep{vaci2019}, nonlinear function approximators \citep{xiao2025} or physical and technical indicators \citep{kalen2021}. These studies describe the average career and treat between-player differences as random deviations from it. A smaller body of work asks whether careers fall into distinct patterns. \citet{wakim2014} applied functional principal component analysis to Win Shares trajectories of NBA players with at least eight seasons and, after $k$-means on the leading scores, distinguished early, middle and late peakers, with the number of groups chosen by a permutation test; \citet{elijah2025} derived player-specific aging curves in ice hockey while modelling the selection that truncates weaker careers; typologies of athlete development have been derived in other sports \citep{boccia2017,gulbin2013}, and group-based trajectory modelling \citep{nagin2005} offers a parametric alternative. All of these use fixed-length or age-aligned representations; none has examined variable-length careers with a learned representation, and none has reported the reproducibility of the resulting typology under resampling. Career length is itself an outcome studied in sports economics: college production and draft position predict NBA careers only partially \citep{coates2010,berri2011draft,motomura2016}, and consistency of performance is rewarded by longer retention \citep{deutscher2017}, which motivates draft status and honours as external validators and cautions against attributing trajectory shape to a single cause.

\paragraph{Time-series clustering and learned representations.} Time series can be clustered with elastic distances such as dynamic time warping \citep{sakoe1978,petitjean2011}, with hand-crafted features, or with representations learned by neural networks \citep{aghabozorgi2015,langkvist2014}. Autoencoders compress a sequence into a code that must retain what reconstruction needs \citep{hinton2006}; Transformer encoders have been used for unsupervised representation learning of series \citep{zerveas2021,vaswani2017}, and dedicated deep clustering models add a clustering objective \citep{ma2019}. The largest comparative study, 300 architectures on 158 datasets, found that deep representations beat classical methods only in some settings \citep{lafabregue2022}, which is why RQ4 compares them on equal terms. Pipelines that pass embeddings through UMAP before HDBSCAN are widespread \citep{angelov2020,grootendorst2022,fragkedaki2024,allaoui2020}; UMAP preserves local neighbourhoods but distorts global distances and can fabricate compact groups at low dimension \citep{chari2023}, HDBSCAN's excess-of-mass extraction can return one large root cluster whose subdivision changes between runs \citep{campello2015,mcinnes2017hdbscan}, and internal indices such as the Silhouette \citep{rousseeuw1987}, Calinski--Harabasz \citep{calinski1974} and Davies--Bouldin \citep{davies1979} assume compact convex clusters; DBCV \citep{moulavi2014} was proposed for density-based solutions.

\paragraph{Reproducibility of cluster solutions.} Stability under perturbation is a long-standing criterion for the number of clusters and the credibility of a clustering \citep{benhur2002,vonluxburg2010}. Consensus clustering accumulates co-assignments over resampled runs and extracts a partition from the co-association matrix \citep{monti2003,fred2005,strehl2002}; \citet{senbabaoglu2014} showed that it can appear to find structure in structureless data and proposed the proportion of ambiguously clustered pairs (PAC) as a diagnostic. We use these tools together with chance-corrected agreement indices \citep{hubert1985,vinh2010}.

\paragraph{Gap.} Prior work provides population aging curves, a few trajectory typologies on fixed-length functional representations, a large toolbox of representation-learning and density-based methods whose fragility is increasingly documented, and reproducibility criteria developed mainly in genomics. Missing is a study that applies a learned representation to complete, variable-length NBA careers and asks, with prespecified criteria, whether the resulting types are reproducible, whether they beat simpler representations, whether they reflect discrete structure at all, and how they relate to external information.
